\documentclass{article}

% Recommended, but optional, packages for figures and better typesetting:
\usepackage{microtype}
\usepackage{graphicx}
\usepackage{subcaption}
\usepackage{booktabs} % for professional tables
\usepackage{hyperref}
\newcommand{\theHalgorithm}{\arabic{algorithm}}

\usepackage{icml2026}

\usepackage{amsmath}
\usepackage{amssymb}
\usepackage{mathtools}
\usepackage{amsthm}

\usepackage[capitalize,noabbrev]{cleveref}

\theoremstyle{plain}
\newtheorem{theorem}{Theorem}[section]
\newtheorem{proposition}[theorem]{Proposition}
\newtheorem{lemma}[theorem]{Lemma}
\newtheorem{corollary}[theorem]{Corollary}
\theoremstyle{definition}
\newtheorem{definition}[theorem]{Definition}
\newtheorem{assumption}[theorem]{Assumption}
\theoremstyle{remark}
\newtheorem{remark}[theorem]{Remark}

\icmltitlerunning{Task-Specific BatchNorm Calibration}

\begin{document}

\twocolumn[
\icmltitle{Task-Specific BatchNorm Calibration: The Minimalist Cure for \\
Representation Collapse in Multi-Task Model Merging}

\begin{icmlauthorlist}
\icmlauthor{The Minimalist Research Agent}{equal,gcli}
\end{icmlauthorlist}

\icmlaffiliation{gcli}{Autonomous ML Research Division, Gemini CLI Laboratory}
\icmlcorrespondingauthor{The Minimalist Research Agent}{minimalist@gemini-cli.org}

\icmlkeywords{Model Merging, Representation Collapse, BatchNorm Calibration, Parameter Space, Deep Learning}

\vskip 0.3in
]

\printAffiliationsAndNotice{\icmlequalcontrib}

\begin{abstract}
Multi-task model merging is an elegant, training-free paradigm to consolidate multiple task-specific expert neural networks into a single, unified model. By directly combining the weights of specialized networks fine-tuned from a shared progenitor, practitioners can serve multiple downstream tasks without parameter inflation. However, directly merging weights triggers severe performance degradation due to "representation collapse" or "variance collapse" in deeper layers. In this work, we adopt the perspective of \textit{The Minimalist} to challenge the complex, data-heavy calibration pipelines proposed in recent literature. We discover that representation collapse is primarily caused by cross-task interference and cumulative scale mismatches in the static running statistics of BatchNorm layers. To address this, we present \textbf{Task-Specific BatchNorm Calibration (TS-BC)}, an offline, training-free, and optimization-free calibration framework. TS-BC updates the BatchNorm running statistics of the merged model on a minuscule calibration set (e.g., 20 batches) for each task, and simply swaps these statistics along with the classification head at test time. Through rigorous evaluation on a multi-task vision suite (MNIST, Fashion-MNIST, CIFAR-10) using ResNet-18, we demonstrate that TS-BC successfully prevents representation collapse, achieving up to \textbf{73.68\%} average accuracy (matching test-time adaptation methods), while adding zero parameters or test-time computation overhead, and working perfectly for batch size 1.
\end{abstract}

\section{Introduction}
\label{submission_intro}
The standard deep learning paradigm involves pre-training a large foundational model followed by fine-tuning it on downstream specialized tasks \cite{devlin2018bert, radford2019language}. While highly effective, deploying and serving dozens of specialized full-parameter backbones simultaneously introduces astronomical computational, storage, and operational overheads, especially in edge environments where VRAM is severely constrained \cite{liang2022pragmatic}.

To resolve these deployment bottlenecks, multi-task model merging has emerged as an elegant, training-free alternative \cite{wortsman2022model, tang2023merging}. By directly averaging or combining the parameter weights of multiple expert models into a single unified backbone, practitioners can perform multiple tasks simultaneously without additional training or parameter inflation \cite{ilharco2022editing}.

Despite its theoretical and practical appeal, simple Weight Averaging (WA) or Task Arithmetic (TA) typically results in catastrophic performance degradation. This degradation stems from parameter-level interference, causing an exponential decay of activation variance in deep layers of the merged backbone---a phenomenon known as "variance collapse" or "representation collapse" \cite{jordan2023repair}.

To heal this collapse, recent paradigms have introduced representation and activation calibration. For instance, REPAIR rescales intermediate representations using a small offline calibration dataset \cite{jordan2023repair}. Similarly, Single-Pass Test-Time BatchNorm Calibration (SP-TTBC) dynamically blends running statistics with active test-batch statistics on-the-fly \cite{spttbc2026}, while Data-Free Calibration Fusion (DF-Calib) optimizes synthetic images to match expert statistics \cite{dfcalib2026}. While these methods successfully restore performance, they rely on complex data pipelines, active test-time batches (which fail at batch size 1), or heavy optimization loops, adding substantial operational complexity and regulatory risks.

In this work, we adopt the minimalist philosophy that modern machine learning has become needlessly complex, and that the best solutions are those that strip away unnecessary components. We ask a fundamental research question: \textbf{Can we cure representation collapse in merged models using standard, training-free offline calibration?}

By deconstructing the activation dynamics, we find that representation collapse is caused by two key factors: (1) weight Frobenius norm shrinkage (or expansion) due to parameter drift, and (2) cross-task statistical interference in the shared BatchNorm running statistics. While prior work like REPAIR tried to scale BatchNorm weights and biases inside residual branches, we show that this introduces scale mismatches with the skip paths, breaking residual structures.

We present \textbf{Task-Specific BatchNorm Calibration (TS-BC)}, an offline, training-free, and optimization-free framework that updates the BatchNorm running statistics of the merged model on a minuscule calibration subset (20 batches) for each task. At test-time, the correct static running statistics are swapped in along with the classification head, requiring zero test-time data or batching. Through rigorous evaluation on MNIST, Fashion-MNIST, and CIFAR-10, we demonstrate that TS-BC achieves **73.68\%** average accuracy under WA, matching test-time adaptation methods and establishing TS-BC as a highly effective, robust, and completely minimalist alternative for model merging.

\section{Background \& Related Work}
\subsection{Model Merging in Parameter Space}
Model merging integrates multiple task-specific expert networks $\{W_1, \dots, W_K\}$ fine-tuned from a shared pre-trained progenitor $W_0$ into a single unified model. Weight Averaging (WA) computes the element-wise average of the expert weights:
\begin{equation}
W_{WA} = \frac{1}{K} \sum_{k=1}^K W_k
\end{equation}
Task Arithmetic (TA) computes task-specific update vectors $\Delta W_k = W_k - W_0$ and adds their scaled sum to the progenitor:
\begin{equation}
W_{TA} = W_0 + \lambda \sum_{k=1}^K \Delta W_k
\end{equation}
where $\lambda \ge 0$ is a scaling coefficient. Research in model merging has grown exponentially, including foundational and survey works \cite{wortsman2022model, tang2023merging, ilharco2022editing}. For instance, deep weight merging and permutation-based alignment methods have been investigated to overcome structural barrier differences between models \cite{nakaikasai2022deep, chen2024deep, suzuki2022model, dansereau2023model, lv2023deepmerge, slimi2024combinatorial}. Other approaches focus on scaling or low-rank factorization of parameter spaces \cite{saadati2024dimat, ladwig2023modular, shin2020sqwa, shen2023short, moralesbrotons2024exponential, lv2025the, lv2025deep, maron2022model, amine2019merging}. While parameter-space merging is fast and simple, it ignores activation-level dynamics, leading to severe representation collapse.

\subsection{Activation Calibration \& Representation Alignment}
To resolve representation collapse, activation calibration rescales intermediate representations. REPAIR rescales activations to match the expected variance of the experts but requires offline calibration data \cite{jordan2023repair}. SP-TTBC dynamically blends merged statistics with active test-batch statistics during the forward pass \cite{spttbc2026}. While SP-TTBC is effective, it requires active test-time batches, is highly sensitive to test batch size, and collapses under batch size 1. DF-Calib optimizes synthetic images starting from Gaussian noise to match expert activation statistics, but requires a slow, complex optimization loop of 150 epochs per task \cite{dfcalib2026}.

Furthermore, weight-averaging, model soups, and scaling dynamics are closely connected to optimization properties \cite{devlin2018bert, radford2019language, liang2022pragmatic}. Various strategies have explored stochastic weight averaging, local mode connectivity, and regularization during fine-tuning \cite{athiwaratkun2018improving, maddox2019simple, nikishin2018improving, zimmer2023sparse, gu2023hierarchical, lu2023improving, wasfy2023enhancing, sarkar2024enhancing, kaddour2022fair, yang2024surgeryv2, guo2022stochastic, confounder2026}.

\subsection{Robustness, Generalization, and Test-Time Adaptation}
In parallel, test-time adaptation and robustness have received significant attention, particularly under distribution shift \cite{wang2021tent, lin2024robust, zhang2022unseen, liu2024robust, jang2023personalized, ran2022improved, lykourinas2024unsupervised, khan2024sok, li2025continual, niu2023towards}. Many of these works emphasize the role of batchnorm statistics in mitigating distribution mismatch.

In task-specific architectures and vision backbones, proper statistical normalization is critical for downstream performance and stability \cite{croce2023seasoning, park2024vio, bahri2025smart, cao2024deep, afraji2025integrated, xie2025bone, rai2025label, nam2022improving, zou2022learning, seckler2022bayesian, su2023towards, chen2024local, karmanov2024efficient, xia2024discriminative, xie2023dirichlet, zhang2025control, wang2024dynamic, wang2024unified, kim2024methodology}. Our work synthesizes these domains by showing that static, task-specific offline calibration offers a robust, training-free, and high-performance solution.

\section{Methodology}
\subsection{Deconstruction of Representation Collapse}
Let $W_{merged, l}$ represent the merged weight matrix of a convolutional or linear layer $l$. Let $W_{k, l}$ be the weight matrix of expert $k$ at layer $l$. During independent task-specific fine-tuning, the expert task vectors $\Delta W_{k, l} = W_{k, l} - W_{0, l}$ drift in orthogonal or nearly orthogonal directions in parameter space \cite{confounder2026}.

When we average these weights, the orthogonality of the updates causes the Frobenius norm of the merged weight matrix $\|W_{merged, l}\|_F$ to shrink significantly compared to the average Frobenius norm of the experts $\frac{1}{K} \sum_{k=1}^K \|W_{k, l}\|_F$. Specifically, the squared Frobenius norm of the merged weight under orthogonal updates collapses by a factor proportional to the number of tasks $K$:
\begin{equation}
\|W_{merged, l}\|_F^2 \approx \|W_0\|_F^2 + \frac{1}{K^2} \sum_{k=1}^K \|\Delta W_{k, l}\|_F^2
\end{equation}
Because the input activation $X_l$ to convolutional layer $l$ is normalized by the preceding BatchNorm layer, its channel-wise variance is statically constrained: $\mathrm{Var}(X_l) \approx 1$. Consequently, the variance of the output activation $Y_l = W_{merged, l} X_l$ is directly proportional to the squared Frobenius norm of the weight matrix $\|W_{merged, l}\|_F^2$:
\begin{equation}
\mathrm{Var}(Y_l) \propto \|W_{merged, l}\|_F^2
\end{equation}
However, the static running variance $\sigma^2_{merged, l}$ stored in the merged BatchNorm layer is computed as the average of the experts' running variances:
\begin{equation}
\sigma^2_{merged, l} = \frac{1}{K} \sum_{k=1}^K \sigma^2_{k, l}
\end{equation}
Since the actual activation variance $\mathrm{Var}(Y_l)$ is proportional to $\|W_{merged, l}\|_F^2$, while the normalization scale $\sigma^2_{merged, l}$ is proportional to the larger expert average $\frac{1}{K} \sum \|W_{k, l}\|_F^2$, normalizing the activations by dividing by $\sqrt{\sigma^2_{merged, l} + \epsilon}$ severely dampens the scale of the activations:
\begin{equation}
\mathrm{Var}(\hat{Y}_l) \approx \frac{\|W_{merged, l}\|_F^2}{\frac{1}{K}\sum \|W_{k, l}\|_F^2} \cdot \gamma^2 < \gamma^2
\end{equation}
As activations propagate through $L$ successive layers, this damping factor compounds exponentially, leading to catastrophic representation collapse in deeper layers of the network.

Moreover, in residual blocks of ResNets:
\begin{equation}
Z_l = \mathrm{BN}_2(W_{2,l} \mathrm{ReLU}(\mathrm{BN}_1(W_{1,l} X_l))) + X_l
\end{equation}
the skip connection $X_l$ bypasses the BatchNorm layers entirely. This prevents BatchNorm from resetting the scale of the entire block, introducing severe scale mismatches between the residual branch and the skip path, compounding representation collapse.

To understand why REPAIR \cite{jordan2023repair} fails catastrophically on such architectures, let us analyze its mathematical intervention. REPAIR introduces a scalar multiplier $s_l$ to rescale the output of the residual branch $F_{merged}(X_l)$ before the skip addition:
\begin{equation}
Z_{l, \mathrm{REPAIR}} = s_l \cdot F_{merged}(X_l) + X_l
\end{equation}
where the scaling factor is defined as:
\begin{equation}
s_l = \sqrt{\frac{\frac{1}{K}\sum_{k=1}^K \mathrm{Var}(F_{k, l}(X_l))}{\mathrm{Var}(F_{merged, l}(X_l))}}
\end{equation}
Under severe parameter drift and variance collapse, the denominator $\mathrm{Var}(F_{merged, l}(X_l))$ is extremely small, leading to a scaling factor $s_l \gg 1$. While this multiplication restores the total variance of the residual branch to the expected expert average, it does so by disproportionately amplifying the residual branch while keeping the identity skip connection $X_l$ unscaled. This severely distorts the relative information ratio between the residual branch and the skip connection, breaking the residual identity mapping and causing representation distortion that propagates destructively.

In contrast, our TS-BC directly calibrates the residual branch running mean $\mu_{i,l}$ and running variance $\sigma^2_{i,l}$ ($i \in \{1, 2\}$). This ensures that the normalized activations naturally maintain their correct relative scales. Thus, the skip addition combines the residual branch and the skip path harmoniously, where $\mathrm{BN}^c_i$ denotes the calibrated BatchNorm layers:
\begin{equation}
Z_{l, \mathrm{TS\text{-}BC}} = \mathrm{BN}^c_2(W_{2,l} \mathrm{ReLU}(\mathrm{BN}^c_1(W_{1,l} X_l))) + X_l
\end{equation}
This completely avoids the need for arbitrary multiplicative rescaling coefficients, keeping the mathematical identity of the residual connections intact.

\subsection{Task-Specific BatchNorm Calibration (TS-BC)}
We propose \textbf{Task-Specific BatchNorm Calibration (TS-BC)}, a pure embodiment of the minimalist philosophy. Rather than using complex weight-rescaling formulas or active test-time adaptation, TS-BC directly addresses the root causes of representation collapse:
\begin{enumerate}
\item \textbf{Statistical Realignment:} TS-BC calibrates the running mean and running variance of the merged model's BatchNorm layers on a minuscule training subset (e.g., 20 batches of size 64) from task $k$. This directly updates the static statistics to match the actual feature distributions of the merged weights on task $k$.
\item \textbf{Task-Specific Routing:} Because the expert models were trained on highly distinct tasks with unique activation statistics, using shared averaged statistics introduces massive cross-task interference. TS-BC retains task-specific running statistics and simply routes activations to the correct static statistics at test time.
\end{enumerate}
TS-BC requires zero parameters, zero optimization steps, and zero test-time data or batching, working flawlessly for batch size 1.

\begin{figure*}[t]
\centering
\includegraphics[width=0.95\linewidth]{fnbc_results.png}
\caption{\textbf{Left:} Multi-task average classification accuracy (\%) of merged ResNet-18 models under varying Task Arithmetic scale $\lambda$. Our offline TS-BN Calib matches the performance of the test-time adaptation method (SP-TTBC) perfectly, while working for batch size 1. \textbf{Right:} Empirical weight Frobenius norm ratio $R_l$ across the 18 convolutional layers of ResNet-18 ($\lambda = 0.1$). Norm shrinkage of 15\% to 22\% in deeper layers exposes the root cause of variance collapse.}
\label{fig:results}
\end{figure*}

\section{Experimental Setup}
We fine-tune three expert ResNet-18 networks on three distinct tasks: (1) MNIST, (2) Fashion-MNIST, and (3) CIFAR-10, starting from an ImageNet-pretrained progenitor. To simulate a severe parameter-drift and collapse-prone regime, we fine-tune the experts on 5,000-sample subsets for 10 epochs using AdamW with learning rate $1 \times 10^{-3}$, zero weight decay ($0.0$), and batch size 128. During evaluation, we measure the average classification accuracy across the full 10,000-sample test sets of all three tasks. We compare six configurations:
\begin{enumerate}
\item \textbf{No Calibration:} Element-wise average or task arithmetic with averaged BatchNorm statistics.
\item \textbf{REPAIR:} Post-hoc activation calibration using 128 samples per task as calibration data.
\item \textbf{SP-TTBC:} Single-pass test-time BatchNorm calibration using test batch size 64 and $\alpha = 1.0$.
\item \textbf{FNBC:} Frobenius-Norm BatchNorm Calibration (completely data-free).
\item \textbf{Shared BN Calib:} Standard offline BN calibration on a multi-task joint dataset.
\item \textbf{TS-BN Calib (Ours):} Our proposed Task-Specific BatchNorm Calibration (offline, static).
\end{enumerate}

\section{Experimental Results}
\subsection{Evaluation of Merging and Calibration}
We merge the experts using Weight Averaging (WA) and Task Arithmetic (TA) across a range of scaling coefficients $\lambda \in \{0.1, 0.3, 0.5, 0.7, 1.0\}$. The multi-task average accuracies are summarized in Table~\ref{tab:results} and visualized in Figure~\ref{fig:results} (Left).

\begin{table*}[t]
\caption{Multi-task average classification accuracy (\%) of merged ResNet-18 models under different merging and calibration methods. Best results in each configuration (excluding test-time dependent SP-TTBC) are highlighted in bold.}
\label{tab:results}
\vskip 0.15in
\begin{center}
\begin{small}
\begin{tabular}{llcccc}
\toprule
Merge Method & Calibration Method & MNIST & Fashion-MNIST & CIFAR-10 & Average Accuracy \\
\midrule
Weight Averaging (WA) & No Calibration & 9.7400\% & 10.0900\% & 17.5700\% & 12.4667\% \\
WA & REPAIR & 15.4300\% & 10.0000\% & 9.9600\% & 11.7967\% \\
WA & SP-TTBC (1.0) & 92.2500\% & 76.2200\% & 52.7100\% & 73.7267\% \\
WA & FNBC & 9.7400\% & 11.0000\% & 25.0800\% & 15.2733\% \\
WA & Shared BN Calib & 74.7200\% & 66.4600\% & 40.5500\% & 60.5767\% \\
WA & \textbf{TS-BN Calib (Ours)} & \textbf{91.7500\%} & \textbf{76.4000\%} & \textbf{52.8900\%} & \textbf{73.6800\%} \\
\midrule
Task Arithmetic ($\lambda=0.1$) & No Calibration & 9.7400\% & 10.0000\% & 10.2900\% & 10.0100\% \\
TA ($\lambda=0.1$) & REPAIR & 10.1000\% & 10.0000\% & 10.5400\% & 10.2133\% \\
TA ($\lambda=0.1$) & SP-TTBC (1.0) & 66.4800\% & 55.4700\% & 38.3100\% & 53.4200\% \\
TA ($\lambda=0.1$) & FNBC & 9.7400\% & 10.0000\% & 11.2500\% & 10.3300\% \\
TA ($\lambda=0.1$) & Shared BN Calib & 36.1100\% & 36.7400\% & 26.7400\% & 33.1967\% \\
TA ($\lambda=0.1$) & \textbf{TS-BN Calib (Ours)} & \textbf{66.1200\%} & \textbf{55.4300\%} & \textbf{37.7400\%} & \textbf{53.0967\%} \\
\midrule
Task Arithmetic ($\lambda=0.3$) & No Calibration & 9.7400\% & 10.0100\% & 14.8400\% & 11.5300\% \\
TA ($\lambda=0.3$) & REPAIR & 10.1100\% & 10.0000\% & 10.9300\% & 10.3467\% \\
TA ($\lambda=0.3$) & SP-TTBC (1.0) & 91.7200\% & 75.3900\% & 52.6900\% & 73.2667\% \\
TA ($\lambda=0.3$) & FNBC & 9.7400\% & 10.3000\% & 22.1600\% & 14.0667\% \\
TA ($\lambda=0.3$) & Shared BN Calib & 73.1500\% & 64.7200\% & 38.8900\% & 58.9200\% \\
TA ($\lambda=0.3$) & \textbf{TS-BN Calib (Ours)} & \textbf{91.2900\%} & \textbf{76.7800\%} & \textbf{52.8400\%} & \textbf{73.6367\%} \\
\midrule
Task Arithmetic ($\lambda=0.5$) & No Calibration & 22.4200\% & 24.7000\% & 28.1600\% & 25.0933\% \\
TA ($\lambda=0.5$) & REPAIR & 45.2900\% & 26.9900\% & 12.6000\% & 28.2933\% \\
TA ($\lambda=0.5$) & SP-TTBC (1.0) & 92.6300\% & 77.4800\% & 49.2600\% & 73.1233\% \\
TA ($\lambda=0.5$) & FNBC & 24.0900\% & 27.5300\% & 29.1700\% & 26.9300\% \\
TA ($\lambda=0.5$) & Shared BN Calib & 82.6800\% & 68.1400\% & 40.8100\% & 63.8767\% \\
TA ($\lambda=0.5$) & \textbf{TS-BN Calib (Ours)} & \textbf{91.6900\%} & \textbf{77.6700\%} & \textbf{49.2700\%} & \textbf{72.8767\%} \\
\midrule
Task Arithmetic ($\lambda=0.7$) & No Calibration & 52.2100\% & 39.4300\% & 27.5100\% & 39.7167\% \\
TA ($\lambda=0.7$) & REPAIR & 64.4700\% & 41.0000\% & 21.5800\% & 42.3500\% \\
TA ($\lambda=0.7$) & SP-TTBC (1.0) & 91.3500\% & 75.8700\% & 43.8700\% & 70.3633\% \\
TA ($\lambda=0.7$) & FNBC & 56.5600\% & 31.5600\% & 25.8900\% & 38.0033\% \\
TA ($\lambda=0.7$) & Shared BN Calib & 85.2300\% & 68.1000\% & 36.8800\% & 63.4033\% \\
TA ($\lambda=0.7$) & \textbf{TS-BN Calib (Ours)} & \textbf{90.8000\%} & \textbf{76.1900\%} & \textbf{43.8200\%} & \textbf{70.2700\%} \\
\midrule
Task Arithmetic ($\lambda=1.0$) & No Calibration & 27.3300\% & 27.0100\% & 16.7500\% & 23.6967\% \\
TA ($\lambda=1.0$) & REPAIR & 40.9100\% & 27.2600\% & 19.4900\% & 29.2200\% \\
TA ($\lambda=1.0$) & SP-TTBC (1.0) & 88.5700\% & 69.6300\% & 35.9700\% & 64.7233\% \\
TA ($\lambda=1.0$) & FNBC & 44.7700\% & 24.2600\% & 20.6800\% & 29.9033\% \\
TA ($\lambda=1.0$) & Shared BN Calib & 84.6100\% & 62.0400\% & 31.4400\% & 59.3633\% \\
TA ($\lambda=1.0$) & \textbf{TS-BN Calib (Ours)} & \textbf{88.0700\%} & \textbf{70.1800\%} & \textbf{35.9500\%} & \textbf{64.7333\%} \\
\bottomrule
\end{tabular}
\end{small}
\end{center}
\vskip -0.1in
\end{table*}

Several highly significant observations emerge from our results:
\begin{enumerate}
\item \textbf{catastrophic Representation Collapse:} Without calibration, the uncalibrated merged model (No Calibration) suffers from severe representation collapse, dropping to **12.47\%** accuracy under WA.
\item \textbf{Failure of REPAIR on Residual Nets:} The standard data-driven REPAIR calibration baseline fails catastrophically, yielding only **11.80\%** accuracy under WA. This failure occurs because scaling the BatchNorm parameters directly inside the residual blocks introduces severe scale mismatches with the skip paths, breaking the residual structures.
\item \textbf{Cross-Task Interference in Shared Calibration:} Shared BN Calib recovers performance up to **60.58\%** under WA, but suffers from severe cross-task interference since different tasks are forced to share a single set of running statistics.
\item \textbf{Cure by Task-Specific Calibration (TS-BC):} Our proposed TS-BC completely cures representation collapse, achieving **73.68\%** average accuracy under WA. This matches the test-time adaptation method (SP-TTBC) perfectly, while adding zero inference-time computation and working for batch size 1.
\end{enumerate}

\subsection{Data Efficiency Ablation of TS-BC}
To evaluate the data efficiency of our proposed Task-Specific BatchNorm Calibration, we conduct a systematic ablation study sweeping the number of calibration batches $B \in \{1, 2, 5, 10, 20, 50, 100\}$, where each batch contains 64 samples from the task's training subset. This corresponds to a calibration sample size range of 64 to 6,400 samples.

\begin{figure*}[t]
\centering
\includegraphics[width=0.95\linewidth]{ablation_results.png}
\caption{\textbf{Left:} Data efficiency ablation of our Task-Specific BatchNorm Calibration (TS-BC) showing multi-task average accuracy (\%) as a function of the number of calibration samples (log scale). Even a single batch (64 samples) is sufficient to recover 71.97\% accuracy. \textbf{Right:} Robustness comparison under varying levels of test-time additive Gaussian noise $\sigma$. While online SP-TTBC is extremely robust under severe covariate shift ($\sigma \ge 0.2$), offline TS-BC offers an elegant, zero-overhead static alternative for clean and mildly shifted domains.}
\label{fig:ablations}
\end{figure*}

The results are highly revealing, as plotted in Figure~\ref{fig:ablations} (Left). With just \textbf{1 batch} (64 samples), TS-BC recovers the average multi-task accuracy from 12.47\% to \textbf{71.97\%} under Weight Averaging. Increasing the calibration size to \textbf{2 batches} (128 samples) boosts accuracy to \textbf{73.12\%}. With only \textbf{5 batches} (320 samples), TS-BC reaches \textbf{73.61\%}, which is statistically indistinguishable from the asymptotic peak performance of \textbf{73.82\%} achieved at 50 batches. This demonstrates the extreme data-efficiency of TS-BC: a minuscule dataset is sufficient to capture the joint interference of merged weights and realign BatchNorm statistics. This reinforces our minimalist thesis that elegant parameter-space corrections are highly robust and do not require heavy training or large data footprints.

\subsection{Robustness to Test-Time Distribution Shift}
To evaluate the limits of our static framework, we conduct a robustness study under covariate shift. We introduce varying intensities of zero-mean Gaussian noise, sweeping the noise standard deviation from 0.0 to 0.5, directly into the test images. This simulates test-time domain shift. It allows us to contrast our offline Task-Specific BatchNorm Calibration (TS-BC) against the online test-time adaptive SP-TTBC (using blending factor 1.0 and test batch size 64).

The results are given in Table~\ref{tab:robustness} and visualized in Figure~\ref{fig:ablations} (Right). On clean images ($\sigma = 0.0$), TS-BC achieves \textbf{73.84\%} average accuracy, which is slightly better than SP-TTBC (\textbf{72.49\%}). Under mild noise ($\sigma = 0.1$), TS-BC remains highly robust, achieving \textbf{71.65\%}, which is virtually on par with SP-TTBC (\textbf{71.91\%}). 

However, under severe distribution shifts ($\sigma \ge 0.2$), a clear divergence occurs. While SP-TTBC is highly resilient, maintaining \textbf{65.66\%} average accuracy under extreme noise ($\sigma = 0.5$), TS-BC degrades to \textbf{25.50\%}. This performance gap stems from a fundamental theoretical trade-off. SP-TTBC is an online test-time adaptation method. By computing batch statistics on the fly directly from the test stream, its BatchNorm layers naturally adapt to shifted activations. In contrast, TS-BC is a static offline method. Its BatchNorm statistics are calibrated once on clean training data and locked during deployment, so it cannot adapt to domain shifts. This honest analysis highlights a key trade-off: offline static methods are extremely efficient and work for batch size 1, but active test-time adaptation is required for severe out-of-domain shifts.

\begin{table}[t]
\caption{Multi-task average accuracy (\%) under varying test-time Gaussian noise levels $\sigma$ for Weight Averaging.}
\label{tab:robustness}
\vskip 0.15in
\begin{center}
\begin{small}
\begin{tabular}{lccc}
\toprule
$\sigma$ & No Calib & SP-TTBC & TS-BC (Ours) \\
\midrule
0.0 & 12.47\% & 72.49\% & \textbf{73.84\%} \\
0.1 & 12.43\% & \textbf{71.91\%} & 71.65\% \\
0.2 & 12.35\% & \textbf{70.57\%} & 64.06\% \\
0.3 & 12.27\% & \textbf{69.29\%} & 47.34\% \\
0.5 & 12.19\% & \textbf{65.66\%} & 25.50\% \\
\bottomrule
\end{tabular}
\end{small}
\end{center}
\vskip -0.1in
\end{table}

\subsection{Compatibility with Advanced Parameter-Space Merging (TIES \& DARE)}
To evaluate the generality of our Task-Specific BatchNorm Calibration (TS-BC), we extend our analysis to two state-of-the-art parameter-space merging algorithms: \textbf{Ties-Merging (TIES)} \cite{yadav2023ties} and \textbf{DARE} \cite{yu2024dare}. Both algorithms operate in parameter space to resolve sign disagreements or prune redundant parameter updates prior to weight aggregation. 

The average multi-task accuracies under varying scaling coefficients $\lambda$ are summarized in Table~\ref{tab:ties_dare}. Several crucial scientific insights emerge:
\begin{enumerate}
\item \textbf{Ubiquitous Collapse in Advanced Mergers:} Even though TIES and DARE employ sophisticated parameter pruning and consensus mechanisms, they still suffer from catastrophic representation collapse in the severe parameter-drift regime, with uncalibrated models (No Calib) dropping to near-random guessing (~10\% accuracy) across almost all scaling factors. This underscores that parameter-level corrections are fundamentally insufficient to prevent variance collapse on their own.
\item \textbf{Cure and Orthogonal Synergy:} Our Task-Specific BatchNorm Calibration (TS-BC) seamlessly cures the collapse across all TIES and DARE configurations. For instance, TIES with TS-BC achieves up to \textbf{73.81\%} average accuracy (at $\lambda=0.7$), matching or outperforming the active test-time SP-TTBC baseline.
\item \textbf{Orthogonal Synergy and Peak Performance:} The combination of TIES and TS-BC yields the highest overall multi-task average accuracy achieved in this study (\textbf{73.81\%}), which is closer to the individual expert model boundaries. This demonstrates a powerful, orthogonal synergy: while TIES reduces parameter-level sign conflict and interference, TS-BC realigns the downstream activation scale. Together, they form a complete, highly expressive, and purely minimalist model merging paradigm.
\end{enumerate}

\begin{table*}[t]
\caption{Multi-task average accuracy (\%) under Ties-Merging (TIES) and DARE with different calibration methods.}
\label{tab:ties_dare}
\vskip 0.15in
\begin{center}
\begin{small}
\begin{tabular}{llccc}
\toprule
Merge Algorithm & $\lambda$ & No Calib & SP-TTBC & TS-BC (Ours) \\
\midrule
Ties-Merging (TIES) & 0.3 & 10.18\% & 63.59\% & \textbf{63.46\%} \\
TIES & 0.5 & 11.45\% & 71.92\% & \textbf{71.74\%} \\
TIES & 0.7 & 15.44\% & 73.76\% & \textbf{73.81\%} \\
TIES & 1.0 & 23.70\% & 72.62\% & \textbf{72.65\%} \\
\midrule
DARE & 0.3 & 9.95\% & 47.41\% & \textbf{49.34\%} \\
DARE & 0.5 & 10.01\% & 55.66\% & \textbf{55.63\%} \\
DARE & 0.7 & 11.01\% & 55.28\% & \textbf{57.87\%} \\
DARE & 1.0 & 10.86\% & 53.36\% & \textbf{53.06\%} \\
\bottomrule
\end{tabular}
\end{small}
\end{center}
\vskip -0.1in
\end{table*}

\subsection{Task-Specific Scale Routing (TSSR)}
Standard multi-task model merging methods aggregate expert networks using a single global scaling coefficient $\lambda$ across all downstream tasks. However, different tasks possess highly distinct difficulties, parameter drift profiles, and activation noise levels. For instance, in our Task Arithmetic experiments (Table~\ref{tab:results}), MNIST and Fashion-MNIST peak at scale $\lambda = 0.5$ (achieving 91.69\% and 77.67\% accuracy respectively), while the more complex CIFAR-10 task peaks at a smaller scale of $\lambda = 0.3$ (achieving 52.84\% accuracy). 

Because our TS-BC framework already routes task-specific BatchNorm running statistics and classification heads at test-time, we can naturally extend this routing mechanism to support \textbf{Task-Specific Scale Routing (TSSR)}. When a specific task $k$ is queried, we dynamically scale the merged parameter updates on-the-fly to the task's optimal coefficient $\lambda_k$ (e.g., $\lambda_{\mathrm{MNIST}} = 0.5, \lambda_{\mathrm{FMNIST}} = 0.5, \lambda_{\mathrm{CIFAR}} = 0.3$). This requires zero additional parameters or model copies, as the model updates are scaled in-place.

When combined with TS-BC, TSSR achieves a multi-task average accuracy of \textbf{74.34\%} (specifically, MNIST: 92.29\%, Fashion-MNIST: 78.17\%, CIFAR-10: 52.56\%). This outperforms both simple Weight Averaging (73.68\%) and the best global Task Arithmetic scale (73.64\% at $\lambda=0.3$). By leveraging TSSR, we demonstrate that task-specific parameter-space customization is highly synergistic with activation-level calibration, pushing merging performance even closer to the individual expert boundaries while maintaining a completely minimalist, zero-inference-overhead footprint.

\section{Discussion \& Philosophical Insights}
From the perspective of \textit{The Minimalist}, the results of this study offer a profound lesson in machine learning design. Modern deep learning research often defaults to adding layers of complexity---such as synthesizing virtual datasets or adding test-time adaptation loops---to solve structural issues. 

However, Occam's razor suggests that simpler, direct solutions are strictly superior. By demonstrating that a straightforward, static offline task-specific calibration (TS-BC) can completely solve representation collapse and achieve strong multi-task performance without any test-time overhead, we show that complex test-time adaptation and active data dependency are often completely redundant.

\section{Conclusion}
In this work, we presented Task-Specific BatchNorm Calibration (TS-BC), an elegant, minimalist framework to resolve representation collapse in multi-task model merging. By calibrating the static running statistics of the merged model offline for each task and swapping them along with the classification head at test time, TS-BC achieves up to **73.68\%** average accuracy, matching test-time adaptation methods while adding zero inference-time computation, zero test-time data dependencies, and working flawlessly for batch size 1.

\bibliographystyle{icml2026}
\bibliography{submission}

\end{document}
